\documentclass[11pt]{article}
\usepackage[a4paper,margin=25mm]{geometry}
\usepackage{amsmath,amssymb,amsthm,hyperref}
\newtheorem{theorem}{Theorem}
\newtheorem{lemma}[theorem]{Lemma}
\newcommand{\Q}{\mathbb Q}
\newcommand{\Z}{\mathbb Z}
\newcommand{\BB}{\operatorname{BB}}
\title{A further reduction of the fixed Busy Beaver index for the 163-variable quartic}
\author{Z diophantine in Q project}
\date{20 November 2026}
\begin{document}\maketitle

\begin{abstract}
The Version~22 Busy Beaver record argument and the v86 sharpening
give a fixed 163-variable quartic using $N=\BB(t_0)$, where
$t_0>2^{826}$. Here we replace the rounded record exponent
$2^{616}$ by its explicit exact upper envelope. We obtain a strictly
smaller fixed index $t_1<2^{826}$, and improve the coefficient
height estimate to $\BB(2^{1010})$. The original elliptic, geometric,
arithmetic and finite-state compiler hypotheses are unchanged.
\end{abstract}

\paragraph{Research provenance and credit.}
This proof refines the quantitative Busy Beaver argument in
\emph{Elliot Glazer's agentic research notes and Version~22 manuscript}
\cite{v22}, developed in his research program with OpenAI Codex
agents. Version~22 specifically credits its new Busy Beaver estimate
and receiving-source checks to an OpenAI Codex agent. The
163-variable scalar-point construction and its original coefficient
tables are Jose Brox's, building on Gabe Conant's
shared-product work with Astra~6.2. The exact record-exponent
compression and coefficient wrapper here are later refinements.
We credit these separate contributions rather than attributing
the entire construction or all source proofs to one person.

\section{Explicit parameter}
Set
\begin{equation}\label{eq:parameters}
 \psi=\prod_{j=2}^{90}(3j+1),\quad
 \kappa=27160+16032\psi,\quad
 U=2^\kappa,\quad
 t_1=2^{20}(16\kappa+2^{186}+1)(2^{186}+3).
\end{equation}
These are ordinary, computable positive integers. Exact arithmetic gives
\begin{equation}\label{eq:bits}
 \operatorname{bitlength}(\psi)=602,\qquad
 \operatorname{bitlength}(\kappa)=616,\qquad
 2^{825}<t_1<2^{826}.
\end{equation}
In particular, $t_1$ is specified without referring to any unknown
height constant or Busy Beaver value.

Let $F_N\in\Z[x_1,\ldots,x_{163}]$ denote the exact
quartic of the v82 physical coefficient tables, with
$C=149^N$. Its only coefficients depending on $C$ are
\[
 [x_{78}^2x_{90}]F_N=-6615898C^2,\quad
 [x_{78}x_{90}]F_N=-298C^2,\quad
 [x_{78}^2]F_N=22201C^4.
\]
The other $2786$ coefficients are the fixed integers printed in those
tables. The family has exactly $2789$ nonzero monomials and content one.

\begin{theorem}\label{thm:main}
Under the inherited arithmetic, weighted-graph height and
register-to-binary compiler assertions of the supplied Version~22
manuscript and the v86 proof, the single polynomial
\[
 \boxed{F_{\BB(t_1)}}
\]
defines $\Z$ Diophantinely in $\Q$. It has total degree four and
$163$ variables, including the input. Its coefficient height is
exactly
\[
 H(F_{\BB(t_1)})=149^{4\BB(t_1)+2}
                 <\BB(2^{1010}).
\]
The rounded index $t=2^{826}$ is also sufficient.
\end{theorem}

\section{No change to the quantitative geometric argument}
The weighted-trace graph with parameter
$(r,T)\in\mathbf P^1\times\mathbf P^1$ and four elliptic points has,
in the worst active torsion stratum,
\[
 \dim V=9,\quad m=\dim V+1=10,\quad
 N_A+1=3^8,\quad N_P+1=4,\quad
 \deg V\le10!3^8.
\]
As checked in v22/v86, its positive-weight actual tuples avoid
all positive-dimensional constant translates, and its height data
satisfy $h(V)<2^{150}$ and
$h_{\mathrm{add}},C_{\mathrm{haut}}<100$.
For von Buhren's quantitative Theorem~3.1 write
\[
\begin{aligned}
\Lambda={}&2^{m^2}(8m\dim V)^{m\dim V}
       (N_A+1)(N_P+1)(\deg V)^{3m},\\
\psi={}&\prod_{j=2}^{m\dim V}(3j+1).
\end{aligned}
\]
The exact worst-case values have
$\operatorname{bitlength}(\Lambda)=2004$ and
$\operatorname{bitlength}(\psi)=602$.
The displayed constants in the published theorem satisfy
\[
\begin{aligned}
\log_2(e_1)&<3+2004\psi,&
\log_2(e_2^2)&<2+8016\psi,\\
\log_2\mu&<2+2004\psi,&
\log_2e_3&<27150+16032\psi.
\end{aligned}
\]
In particular, by \eqref{eq:parameters},
\begin{equation}\label{eq:margin}
 e_2^2,e_3,\mu < 2^{\kappa-10}=U/1024<U/100.
\end{equation}
Cover the direction cube $[-1,1]^4$ by cells of diameter at most
$1/e_1$. The same proof as v86 allows at most nine records per cell
and each of at most $256$ torsion strata, so the logarithm of their
total number is less than
\[
 36+8016\psi<27160+16032\psi=\kappa.
\]
Consequently the total record count is strictly less than
$U=2^\kappa$; the original v86 choice $U=2^{2^{616}}$ is unnecessary.

\section{Explicit register and binary machine bounds}
The source selector of Version~22 enumerates finite rational weights
and elliptic point codes in exact rational/quadratic-field arithmetic.
It maintains the record criterion
\[
 S>UR,\qquad S>U(b(r)+b(T)),
\]
and at its actual final record value $R_f$ runs for more than
\[
 G(R_f)=U^3(R_f+3)
\]
steps. By the unchanged v86 arithmetic lemma, $G(R_f)$ is a
sufficient threshold for every larger $N$.
The new $U$ still exceeds all numerical contact constants, satisfies
\eqref{eq:margin}, and is certainly greater than $2^{16}$;
therefore every inequality in that lemma remains applicable.

The source program is changed only by replacing the literal
\texttt{power(2,power(2,1000))} with
\texttt{power(2,kappa)}, where the fixed integer $\kappa$ is
given by \eqref{eq:parameters}, and replacing
\texttt{power(scale,1000)} with \texttt{power(scale,3)}.
Since $\kappa$ has only $616$ bits, its literal initialization fits
within the already established $2^{186}$-instruction fixed-core
budget. The function call graph and tuple structure do not increase.

The ordinal of the actual last record is some $j_*<2^\kappa$.
Hardwiring it requires at most $\kappa$ binary bits.
The two-scratch-counter loader in v22 uses fewer than 16 instructions
per bit; the revised fixed program therefore has
\[
 L<16\kappa+2^{186}\quad\text{counter instructions},
 \qquad R<2^{186}+2\quad\text{counters}.
\]
The source's unary-counter-to-binary-tape compiler consequently
uses strictly fewer than
\[
 2^{20}(L+1)(R+1)
 <2^{20}(16\kappa+2^{186}+1)(2^{186}+3)=t_1
\]
binary working states. The program for the actual (unknown but finite)
ordinal $j_*$ halts after more than $G(R_f)$ transitions.
Padding unused states proves
\[
 G(R_f)<\BB(t_1).
\]
The all-larger-$N$ arithmetic theorem from v86 and the original
ordered integer-completeness proof then give
$F_{\BB(t_1)}(\Q^{163})$ precisely the desired input projection.

\section{Sharper coefficient-height wrapper}
Since $2^{825}<t_1<2^{826}$, a maximizing
$t_1$-state halting machine has $2t_1$ transition entries,
each of width $826+3=829<2^{10}$ bits.
Its hardwired table contains strictly fewer than $2^{837}$ bits.
The source's coefficient-interpreter core remains smaller than
$2^{148}$ instructions; its reusable bit loader consumes fewer
than 16 instructions per bit. Thus
\[
 L_{\rm coeff}<16(2t_1\cdot829)+2^{148}<2^{841},\qquad
 R_{\rm coeff}<2^{149}.
\]
The binary simulation uses strictly fewer than
\[
 2^{20}(L_{\rm coeff}+1)(R_{\rm coeff}+1)
 \le2^{20+841+149}=2^{1010}
\]
working states. This program simulates a maximizing
$t_1$-state halting table, computes
$149^{4\BB(t_1)+2}+1$, executes more than that many
additional transitions and halts. Padding proves
\[
 149^{4\BB(t_1)+2}+1<\BB(2^{1010}).
\]
Because $\BB(t_1)\ge t_1>53$, the exact original
coefficient-height formula
$\max(126616122944180358034867798936,149^{4N+2})$
reduces to $149^{4N+2}$. This completes the theorem.

\paragraph{Scope.}
The explicit constants and new integer comparisons have independent
reproducible checks. The complete analytic geometry, source
elliptic-point realization and explicit register compilation
are inherited from v22/v86, not reproved by those finite checks.

\begin{thebibliography}{3}
\bibitem{v22} Elliot Glazer and attributed project contributors (agentic notes with OpenAI Codex), \emph{Z is Diophantine in Q}, Version~22,
user-supplied \texttt{quartic\_v22.tex}, pinned private commit
\texttt{4b43ba288b9b8639cc5cc85f8c1507d64c698659}.
\bibitem{v86} \emph{A fixed 163-variable quartic with a sharper
explicit Busy Beaver index}, project v86, 20 November 2026,
\texttt{papers/quartic/bb163\_fixed\_20261120\_v86.tex}.
\bibitem{VB} J. von Buhren,
\emph{Borne de hauteur semi-effective pour le
probl\`eme de Mordell--Lang dans une vari\'et\'e ab\'elienne},
J. Th\'eorie des Nombres de Bordeaux \textbf{29} (2017),
307--320, Theorem~3.1, p.~312, DOI: 10.5802/jtnb.981.
\end{thebibliography}
\end{document}
